% Self-contained mathematical section. Requires amsmath, amssymb, amsthm.
% Bibliography keys: ProveItArticle, GRHLean, Liehr2024, Wellershoff2024.
\section{One positive prime weight determines all labelled Dirichlet zeros}
\label{rw:sec:main}

The inverse-problem section of the supplied article constructs a complex
prime weight whose autocorrelation determines all labelled primitive
Dirichlet zero multisets. It asks whether one real nonnegative weight can
do the same \cite[section ``Recovering zeros and characters from prime
data'', final question ``Fewer filters and weights'']{ProveItArticle}.
The answer below is affirmative. Second-order correlation already suffices.
The weight is computable independently of the zeros and stays uniformly
close to one. The construction uses three bounded frequency shifts to
recover the orientation that taking real parts would otherwise erase.

Throughout this section, Dirichlet GRH means that the nontrivial zeros of
every primitive Dirichlet $L$-function lie on $\Re s=1/2$. This is exactly
the input supplied by the main exported theorems of \cite{GRHLean}; no
hypothesis about Artin or automorphic $L$-functions is used. Put
\[
 m_\chi(t)=\operatorname{ord}_{s=1/2+it}L(s,\chi),\qquad
 D(t)=(\tfrac12+it)(\tfrac32+it),\qquad
 K(t)=|D(t)|.
\]
At a nonzero value of $L$, the order is zero. Write $m_\zeta$ for the
conductor-one character. Both signs of $t$ are retained for a nonreal
character.

\begin{theorem}[A real positive universal arithmetic encoder]
\label{rw:thm:encoder}
There is an explicitly specified computable function
$W:\mathbb N_{\ge1}\to[3/4,5/4]$, independent of the zero multisets, such that the
following holds under Dirichlet GRH. Define, for real $t$,
\[
 F(t)=e^{-t/2}\left\{
 \sum_{1\le n\le e^t}\Lambda(n)W(n)(1-ne^{-t})-\frac{e^t}{2}\right\}.
\]
The real autocorrelation
\[
 C(h)=\lim_{T\to\infty}\frac1T\int_0^T F(t+h)F(t)\,dt
 \qquad(h\in\mathbb R)
\]
exists, is continuous, even and positive definite, and determines
$m_\chi(t)$ for every primitive character $\chi$ and every real $t$.
In particular, the labels, multiplicities and signs of the ordinates are
recovered. No simplicity, common-zero exclusion or linear independence
assumption is required.

The same conclusion holds for a sum over primes alone after adding the
explicit bounded correction in Corollary~\ref{rw:cor:primes}.
More generally, for every prescribed rational $0<\varepsilon<1$ one
can require $1-\varepsilon\le W(n)\le1+\varepsilon$ for all $n$.
The determination uses exact limiting correlation data. It does not
assert a uniformly stable reconstruction from a fixed finite observation.
\end{theorem}

\subsection{A separated positive code}

Enumerate all pairs $(\chi_j,a_j)$, where $\chi_j$ is primitive and
nonprincipal and $a_j\in\{0,1,\sqrt2\}$, as $j=1,2,\ldots$, in a fixed
computable order. For example, order the characters by conductor and then
by their finite tables of algebraic values, and list the three shifts
consecutively. Let $q_j$ be the conductor and $w_j=q_j+1$. Put $c_0=1$,
$w_0=2$, $k_0=0$, and let $p_j$ denote the $j$th rational prime.
The construction below is the multiquadratic prefix code of
\cite[lemma ``Prefix separation'']{ProveItArticle}, with a smaller tail
budget needed for a real positive weight. It is included in full.

Inductively set
\begin{align}
 D_n^*&=2^{\max_{0\le j\le n}k_j},&
 d_n&=(D_n^*)^{-1}
       \bigl(3D_n^*(n+1)\bigr)^{-(2^n-1)},\label{rw:eq:dn}\\
 c_j&=2^{-k_j}\sqrt{p_j},\label{rw:eq:cj}
\end{align}
where $k_j$ is the least nonnegative integer satisfying
\begin{equation}\label{rw:eq:codechoice}
 2^{-k_j}p_jw_j\le
 2^{-j-3}\min\left(1,\min_{0\le n<j}\frac{d_n}{n+1}\right).
\end{equation}
Every comparison in \eqref{rw:eq:codechoice} is rational, so the choice is
effective. In particular $d_0=1$, and
\begin{equation}\label{rw:eq:tails}
 \sum_{j\ge1}c_jw_j\le\frac18,\qquad
 \sum_{j>n}c_jw_j\le
 2^{-n-3}\frac{d_n}{n+1}.
\end{equation}

\begin{lemma}[Prefix decoding]\label{rw:lem:code}
If $u_0,\ldots,u_n$ are integers, not all zero, with
$|u_j|\le(n+1)w_j$, then
\[
 \left|\sum_{j=0}^n c_ju_j\right|\ge d_n.
\]
Consequently $E(m)=\sum_{j\ge0}c_jm_j$ is injective on nonnegative integer
sequences satisfying $m_j\le H w_j$ for some finite $H$.
If $n+1\ge H$ and $|A'-E(m)|<d_n/8$, the true prefix is the unique vector
of integers $0\le v_j\le(n+1)w_j$, $0\le j\le n$, satisfying
\[
 \left|A'-\sum_{j=0}^n c_jv_j\right|<\frac{d_n}{2}.
\]
\end{lemma}

\begin{proof}
The field $\mathbb Q(\sqrt{p_1},\ldots,\sqrt{p_n})$ has degree $2^n$;
its automorphisms change the signs of the square roots independently.
The square classes of the distinct primes are independent, as follows
by taking their valuations modulo two. Thus
$z=D_n^*\sum_{j=0}^nc_ju_j$ is a nonzero algebraic integer. Every conjugate
of $z$ has modulus at most
\[
 D_n^*(n+1)\left(w_0+\sum_{j=1}^nc_jw_j\right)
 \le\frac{17}{8}D_n^*(n+1)<3D_n^*(n+1).
\]
The product of its $2^n$ conjugates is a nonzero integer. Bounding all
but the displayed conjugate gives precisely the lower bound $d_n$.
By \eqref{rw:eq:tails}, the omitted tail of $E(m)$ is at most $d_n/8$
when $n+1\ge H$. The true prefix is therefore within $d_n/4$ of $A'$;
every other admissible prefix is more than $3d_n/4$ away. Finally, if two
sequences differed, a sufficiently long prefix containing a difference
would contradict this uniqueness with $A'=E(m)$.
\end{proof}

Define the promised weight by
\begin{equation}\label{rw:eq:weight}
 W(n)=1+2\sum_{j\ge1}c_j
       \Re\bigl(\chi_j(n)n^{-ia_j}\bigr).
\end{equation}
The series converges uniformly in the positive integer $n$.
Equation~\eqref{rw:eq:tails} gives $|W(n)-1|\le2\sum c_j\le1/4$.
It also gives an effective uniform truncation error, proving the stated
computability of $W$. No zero or zero approximation enters this definition.
For the arbitrary-$\varepsilon$ version, replace $2^{-j-3}$ in
\eqref{rw:eq:codechoice} by $\sigma2^{-j}$, where
$\sigma=\min(\varepsilon/2,1/8)$. All separation and tail estimates remain
valid, while $|W(n)-1|\le2\sigma\le\varepsilon$.

\subsection{The shifted explicit formula and the infinite sum}

\begin{lemma}[Uniformly bounded shifts]\label{rw:lem:explicit}
Assume Dirichlet GRH. Let $\chi$ be primitive and nonprincipal of conductor
$q$, and let $|a|\le\sqrt2$. Uniformly for $x\ge1$,
\begin{equation}\label{rw:eq:shifted}
 \sum_{1\le n\le x}\Lambda(n)\chi(n)n^{-ia}(1-n/x)
 =-\sum_\gamma
 \frac{x^{1/2+i(\gamma-a)}}{D(\gamma-a)}
 +O\bigl(x^{1/4}\log(q+2)\bigr).
\end{equation}
The zeros in this formula are repeated with multiplicity. The implied
constant is absolute, and
\begin{equation}\label{rw:eq:abssum}
 \sum_\gamma |D(\gamma-a)|^{-1}\ll\log(q+2)
\end{equation}
with an absolute constant. For $\zeta$ and $a=0$, the same formula holds
with the additional main term $x/2$.
\end{lemma}

\begin{proof}
We give the uniformity details because they justify using infinitely many
characters. The standard completed function is
\[
 \Xi_\chi(s)=\bigl(s(s-1)\bigr)^{\mathbf1_{q=1}}
 (q/\pi)^{(s+\kappa)/2}\Gamma((s+\kappa)/2)L(s,\chi),
 \qquad \chi(-1)=(-1)^\kappa.
\]
Its Hadamard product and functional equation, with the zeros on the
critical line, give
\begin{equation}\label{rw:eq:poisson}
 \Re\frac{\Xi_\chi'}{\Xi_\chi}(2+iu)
 =\sum_\gamma\frac{3/2}{(3/2)^2+(u-\gamma)^2}.
\end{equation}
For example, subtract the logarithmic derivative at $1/2+iu$, whose real
part is zero away from zeros, and then extend by continuity in $u$.
The Euler series at $2+iu$ and the digamma factor show that the left side
is $O(\log(q(|u|+3)))$, absolutely in $q,u$. Thus each interval
$[u-1,u+1]$ contains $O(\log(q(|u|+3)))$ zeros, with multiplicity.
Summing these estimates over unit intervals proves
\eqref{rw:eq:abssum}, uniformly for the indicated bounded shifts.

For completeness a crude numerical multiplicity bound follows directly.
The digamma series implies
$\Re\psi(x+iy)\le2+\log(|y|+2)$ for $1\le x\le3/2$:
split its series after $N=\lceil\max(2,|y|)\rceil$ terms and bound the
remaining terms by $\tfrac12(n+1)^{-2}+y^2(n+1)^{-3}$.
Also $\sum_{n\ge2}\Lambda(n)n^{-2}<2$ and the two possible rational
terms in $\Xi_\chi'/\Xi_\chi$ have real part at most $3/2$ together.
Hence the left side of \eqref{rw:eq:poisson} is at most
$\tfrac92+\tfrac12\log(q(|u|+3))$. Each zero in $[u-1,u+1]$ contributes
at least $6/13$, so the safe bound
\begin{equation}\label{rw:eq:localcount}
 \#\{\gamma:|\gamma-u|\le1\}
 \le10+2\log(q(|u|+3))
\end{equation}
is valid, including $q=1$.

For $x\ge1$ the Mellin inversion formula is
\[
 \frac1{2\pi i}\int_{(2)}
 -\frac{L'}{L}(s+ia,\chi)\frac{x^s}{s(s+1)}\,ds.
\]
Move the line to $\Re s=1/4$. The logarithmic derivative there is
$O(\log(q(|\Im s|+4)))$, uniformly for $|a|\le\sqrt2$.
To see this bound without proximity to zeros, compare the Hadamard
logarithmic derivative at $1/4+iv$ with that at $2+iv$. The difference
is an absolutely convergent sum of terms bounded by
$O((1+|v-\gamma|^2)^{-1})$; the local count just proved and the elementary
gamma terms give the assertion. Horizontal edges can be chosen at
distance $\gg1/\log(q(T+3))$ from zero ordinates, using the same local
count. Their contributions tend to zero; the resulting logarithmic
derivative bound on these edges is $O(\log^2(q(T+3)))$.
The new vertical integral has absolute value at most
\[
 Cx^{1/4}\int_{\mathbb R}
 \frac{\log(q(|v|+4))}{1+v^2}\,dv
 \ll x^{1/4}\log(q+2).
\]
The crossed poles are at $s=1/2+i(\gamma-a)$, and their residues give
\eqref{rw:eq:shifted}. Nonprincipal characters introduce no pole at
$1-ia$. The poles of the kernel and the trivial zeros have real part
at most zero, so none was crossed. For $\zeta$ with $a=0$, its pole
at $s=1$ contributes $x/2$.
\end{proof}

For each real $\lambda$ put
\begin{equation}\label{rw:eq:amplitude}
 A(\lambda)=m_\zeta(\lambda)+
 \sum_{j\ge1}c_j
 \bigl\{m_{\chi_j}(a_j+\lambda)+m_{\chi_j}(a_j-\lambda)\bigr\}.
\end{equation}
These are nonnegative, and $A(-\lambda)=A(\lambda)$.
The conjugate identity $m_{\bar\chi}(v)=m_\chi(-v)$ is what makes the
second contribution of $2\Re(\chi(n)n^{-ia})$ equal to
$m_\chi(a-\lambda)$.
Since $\sum c_j\log(q_j+2)<\infty$, Lemma~\ref{rw:lem:explicit} and
absolute convergence yield
\begin{equation}\label{rw:eq:Fourier}
 F(t)=G(t)+O(e^{-t/4}),\qquad
 G(t)=-\sum_{\lambda}\frac{A(\lambda)}{D(\lambda)}e^{i\lambda t},
 \qquad
 \sum_\lambda\frac{A(\lambda)}{|D(\lambda)|}<\infty.
\end{equation}
The sum is over the countable union of the shifted spectra and the zeta
spectrum. That union need not be locally finite. Its absolute convergence
is the justification for all regrouping and averaging here; local
finiteness of the union is not assumed.

Expand the absolutely convergent double series for $G(t+h)G(t)$, using
$G$ real. The mean of $e^{i\omega t}$ is zero unless $\omega=0$.
The decaying error in \eqref{rw:eq:Fourier} does not change the mean.
Consequently
\begin{equation}\label{rw:eq:correlation}
 C(h)=\sum_\lambda\frac{A(\lambda)^2}{K(\lambda)^2}e^{i\lambda h}.
\end{equation}
This is an absolutely convergent series with nonnegative Fourier
coefficients, proving the existence and the stated properties of $C$.
Its Fourier--Bohr coefficients are
\begin{equation}\label{rw:eq:decodeA}
 P(\lambda):=\lim_{U\to\infty}\frac1U\int_0^U
 C(h)e^{-i\lambda h}\,dh
 =\frac{A(\lambda)^2}{K(\lambda)^2},\qquad
 A(\lambda)=K(\lambda)\sqrt{P(\lambda)}.
\end{equation}
Here the positive square root is justified by the nonnegative amplitudes.
At $\lambda=0$ the same formula holds without an extra factor of two.

Equation~\eqref{rw:eq:localcount} implies that the integer vector
\[
 \bigl(m_\zeta(\lambda),
 m_{\chi_j}(a_j+\lambda)+m_{\chi_j}(a_j-\lambda)\ (j\ge1)\bigr)
\]
satisfies the bound in Lemma~\ref{rw:lem:code}. One explicit sufficient
choice is
\begin{equation}\label{rw:eq:Hlambda}
 H_\lambda=\left\lceil30+10\log(|\lambda|+5)\right\rceil.
\end{equation}
Indeed each paired multiplicity is at most
$20+4\log(q_j(|\lambda|+5))\le H_\lambda(q_j+1)$, using
$\log q\le q-1$; the zeta coordinate is smaller than $2H_\lambda$.
The prefix code therefore recovers, for every character, the three
functions
\begin{equation}\label{rw:eq:Sa}
 S_{\chi,a}(x)=m_\chi(a+x)+m_\chi(a-x),
 \qquad a=0,1,\sqrt2.
\end{equation}

\subsection{Recovering an oriented discrete spectrum from reflections}

The reflection principle used next has close precedents in holomorphic
phase retrieval. Liehr characterizes parallel-line uniqueness by exclusion
of an arithmetic progression \cite{Liehr2024}, and Wellershoff treats
three parallel lines with rationally independent distances
\cite{Wellershoff2024}. We include the elementary discrete-measure
argument and an exact stopping rule; the reflection principle itself is
not asserted to be new.

\begin{proposition}[The reflection-group criterion]
\label{rw:prop:reflection}
Let $\mathcal A\subset\mathbb R$ be nonempty and let $a_0\in\mathcal A$.
For locally finite positive atomic measures, the family of centered
symmetrizations about $a\in\mathcal A$ is injective if and only if
\[
 \Gamma=\left\langle2(a-a_0):a\in\mathcal A\right\rangle_{\mathbb Z}
\]
is nondiscrete. In particular, three centers $a_0,a_1,a_2$ with
$(a_1-a_0)/(a_2-a_0)$ irrational suffice, whereas two centers never
suffice for all locally finite positive atomic measures.
\end{proposition}

\begin{proof}
If two measures have the same symmetrizations, their signed difference
$\nu$ satisfies $R_a\nu=-\nu$, where $R_a$ is reflection about $a$.
Composing two reflections gives invariance under every translation in
$\Gamma$. A nondiscrete additive subgroup of $\mathbb R$ is dense.
If $\nu$ had a nonzero atom, translation would give equal nonzero atoms
at a dense set of translates. Infinitely many would lie in one compact
interval, contradicting local finiteness of its total variation.
Since $\nu$ is atomic, it is zero.

Conversely, a discrete subgroup is either $b\mathbb Z$ for $b>0$, or
$\{0\}$. In the first case choose $0<\delta<b/2$ and set
\[
 \mu=\sum_{k\in\mathbb Z}\delta_{a_0+\delta+kb},\qquad
 \widetilde\mu=\sum_{k\in\mathbb Z}\delta_{a_0-\delta+kb}.
\]
These are distinct and every reflection $R_a$, $a\in\mathcal A$,
interchanges them because $2(a-a_0)\in b\mathbb Z$.
Their symmetrizations therefore agree. For $\Gamma=\{0\}$ use the
two single atoms at $a_0\pm\delta$. The lattice counterexamples have
bounded multiplicities and finite $(1+t^2)^{-1}$-weighted mass. They
are abstract spectra, not asserted Dirichlet zero multisets.
\end{proof}

For Dirichlet zeros one can recover each multiplicity by a terminating
finite calculation once exact values of the symmetrizations are supplied.
Let $m:\mathbb R\to\mathbb Z_{\ge0}$ have locally finite support and set
$S_a(x)=m(a+x)+m(a-x)$. Define
\begin{equation}\label{rw:eq:Da}
 \Delta_a(y)=S_a(y+a)-S_0(y)=m(y+2a)-m(y).
\end{equation}
Write
\[
 (1+\sqrt2)^k=p_k+q_k\sqrt2,\qquad
 u_k=2(q_k\sqrt2-p_k)=-2(1-\sqrt2)^k.
\]
Then $p_k,q_k$ are positive integers, $u_k\ne0$, and $u_k\to0$.
The path with $q_k$ steps $+2\sqrt2$ and then $p_k$ steps $-2$ gives
the fully explicit identity
\begin{align}
 T_k(t)
 &:={\sum_{r=0}^{q_k-1}}\Delta_{\sqrt2}(t+2r\sqrt2)
 -{\sum_{s=1}^{p_k}}\Delta_1(t+2q_k\sqrt2-2s)\notag\\
 &=m(t+u_k)-m(t).\label{rw:eq:Pelltelescope}
\end{align}

\begin{proposition}[An exact stopping rule]\label{rw:prop:stop}
For fixed real $t$, inspect $k=1,2,\ldots$ until
$S_0(t+u_k)=0$. This happens for every sufficiently large $k$ and certifies
\[
 m(t)=-T_k(t).
\]
\end{proposition}

\begin{proof}
Local finiteness gives punctured neighborhoods of both $t$ and $-t$
containing no support point. Thus, for all sufficiently large $k$,
both $m(t+u_k)$ and $m(-t-u_k)$ vanish. Conversely, a zero value of
$S_0(t+u_k)$ certifies $m(t+u_k)=0$ because both summands are
nonnegative. Equation~\eqref{rw:eq:Pelltelescope} then proves the formula.
\end{proof}

This stopping rule is an exact-oracle statement; it does not decide from
a noisy real approximation whether a spectral coefficient vanishes.
Applying it to the three functions in \eqref{rw:eq:Sa} finishes the proof
of Theorem~\ref{rw:thm:encoder}. Only the individual character's spectrum
is required to be locally finite.

\subsection{Using primes alone and the scope of the conclusion}

\begin{corollary}[A literal prime statistic]\label{rw:cor:primes}
For the weight \eqref{rw:eq:weight} put
\begin{equation}\label{rw:eq:bias}
 B(t)=\frac23+4\sum_{j:\chi_j^2\ {
m principal}}
 c_j\Re\left\{
 \frac{e^{-ia_jt}}{(1-2ia_j)(3-2ia_j)}\right\}.
\end{equation}
Then the real signal
\[
 F_{\rm p}(t)=e^{-t/2}\left\{
 \sum_{p\le e^t}(\log p)W(p)(1-pe^{-t})-\frac{e^t}{2}\right\}+B(t)
\]
has the autocorrelation $C$ of Theorem~\ref{rw:thm:encoder}.
The correction $B$ is explicit and computable; it is a linear combination
of the constant function and sines and cosines at frequencies $1$ and
$\sqrt2$.
\end{corollary}

\begin{proof}
The terms of exponent at least three in the prime-power sum contribute
$O(x^{1/3}\log^2x)=o(\sqrt x)$, since $W$ is bounded. For a fixed
character and fixed $a$, the prime number theorem for $\chi^2$ and
partial summation give
\[
 x^{-1/2}\sum_{p\le\sqrt x}(\log p)\chi(p)^2p^{-2ia}(1-p^2/x)
 =\mathbf1_{\chi^2\ {
m principal}}
 \frac{2x^{-ia}}{(1-2ia)(3-2ia)}+o(1).
\]
The coefficient is the integral of $u^{-2ia}(1-u^2/x)$ from zero to
$\sqrt x$. The unshifted baseline gives $2/3$. Chebyshev's estimate
bounds every normalized sum absolutely by one common constant,
independent of the character. Since $\sum c_j<\infty$, dominated
convergence permits summing the individual limits. Therefore the
normalized square contribution is $B(\log x)+o(1)$, giving
$F_{\rm p}(t)=F(t)+o(1)$. Both signals are bounded for $t\ge0$, so their
fixed-lag correlations agree.
\end{proof}

The contribution relative to the supplied article is the positive real
arithmetic encoding \eqref{rw:eq:weight} together with recovery of the
oriented labelled spectra from its real autocorrelation. The existing
complex-weight encoder and its multiquadratic separation mechanism are
explicit antecedents. The reflection argument has the published
antecedents just cited. A targeted literature search did not locate this
combined arithmetic theorem; that is evidence for a candidate new result,
not a claim that an exhaustive priority search has been completed.

The coefficient code is deliberately very small and its decoding
thresholds deteriorate rapidly with the character index. The theorem
answers an exact information question. Finite-time spectral estimation
with controlled error, and an economical well-conditioned encoding, are
separate problems. No assertion here proves simplicity or nonvanishing
at the central point, nor does it enlarge Dirichlet GRH to other families.
